\section{Monastic Conversion and the Embassy to Constantinople (c.~574--586)}

\subsection{The long resistance and the Caelian cloister}

Gregory's turn to monastic life was neither a sudden lightning bolt nor a flight
from youthful dissipation. In the Epistle Dedicatory to Leander of Seville
prefixed to the \work{Moralia in Iob}, he left a remarkably candid portrait of
how slow and tangled his conversion had been:

\begin{quote}
For late and long I declined the grace of conversion, and after that I had been
inspired with an heavenly affection I thought it better to be still shrouded in
the secular habit. For though I had now disclosed to me what I should seek of
the love of things eternal, yet long-established custom had so cast its chains
upon me, that I could not change my outward habit: and while my purpose still
compelled me to engage in the service of this world as it were in semblance
only, many influences began to spring up against me from caring for this same
world, so that the tie which kept me to it was now no longer in semblance only,
but what is more serious, in my own mind. At length being anxious to avoid all
these inconveniences, I sought the haven of the monastery, and having left all
that is of the world, as at that time I vainly believed, I came out naked from
the shipwreck of human life. (\work{Moralia in Iob}, Ep.\ Ded.\ 1)
\end{quote}

Few ancient autobiographies diagnose spiritual delay with greater precision.
Gregory had tried to serve God inwardly while keeping the outward habits of
senatorial office, only to discover that external entanglements quietly invade
the interior life. Following his father's death (around 574), he broke the
chain decisively. As Gregory of Tours attests from contemporary Roman report,
he used his inherited wealth to endow six monasteries on his family estates in
Sicily and established a seventh within the walls of Rome, selling the rest of
his property and even his household furnishings to distribute to the poor
(\work{Historia Francorum} X.1). Gregory himself identifies the Roman
foundation in a sermon as ``my monastery, situated beside the Church of the
blessed martyrs John and Paul'' (\work{Homiliae in Evangelia} 38.16)---the
monastery of Saint Andrew \latin{ad clivum Scauri} on the Caelian Hill, built
within his own ancestral home.

Inside Saint Andrew's, Gregory did not first claim the authority of abbot; he
entered as a monk under obedience and threw himself into severe asceticism.
Gregory of Tours writes that ``such was his abstinence in eating, such his
watchfulness in prayer, such his zeal in fasting, that his weakened stomach
scarce availed to support his frame'' (\work{Historia Francorum} X.1). Gregory
confirms the lasting toll on his body in the \work{Moralia} preface: ``For many
a year's circuit has gone by since I have been afflicted with frequent pains in
the bowels, and the powers of my stomach being broken down, makes me at all
times and seasons weakly; and under the influence of fevers, slow, but in
constant succession, I draw my breath with difficulty'' (\work{Moralia}, Ep.\
Ded.\ 5). Barmby notes from the \work{Dialogi} (III.33) that Gregory even
recalled fainting during Holy Week from a refusal to break his fast before
Easter until rescued through the prayer of a holy monk named Eleutherius. Whether rooted in constitution, ascetic excess, or both,
chronic gastrointestinal illness and gout shadowed every remaining decade of
his life.

\subsection{Deacon and \emph{apocrisiarius} in Constantinople}

The monastic peace Gregory loved lasted only a few years. Summoned out of his
cloister under obedience, he was ordained one of the seven regional deacons of
Rome and dispatched to Constantinople as papal \latin{apocrisiarius} (resident
envoy, or \latin{responsalis}) at the imperial court. Later biographers diverge
slightly on which pope initiated the call: combining Paul the Deacon and John
the Deacon, Barmby observes that Pope Benedict I (575--579) probably ordained
him deacon, while Pelagius II (579--590) actually sent him east around 578 or
579, where he remained under Emperors Tiberius II (578--582) and Maurice
(582--602) until about 585 or 586 (Barmby, Prolegomena). Gregory himself felt
the call as a violent uprooting: ``as the vessel that is negligently moored is
very often (when the storm waxes violent) tossed by the water out of its
shelter on the safest shore, so under the cloak of the Ecclesiastical office, I
found myself plunged on a sudden in a sea of secular matters''
(\work{Moralia}, Ep.\ Ded.\ 1).

His diplomatic mission was grim. Ravenna's exarch could not or would not
shield central Italy from the Lombards, and Pelagius II urgently instructed
Gregory in 584 to press the imperial court for troops and relief (Barmby,
Prolegomena). Little military aid was forthcoming from an empire hard-pressed
on the Persian and Balkan frontiers. Yet the six-year residency in the eastern
capital proved decisive for Gregory's future. He learned the machinery of the
imperial palace and formed lasting friendships with figures at the summit of
Byzantine society: Emperor Maurice himself---whose son Gregory held at the
baptismal font, making them spiritual kinsmen (\work{Historia Francorum}
X.1)---Maurice's devout sister Theoctista, the court physician Theodorus, the
patrician Narses, and noblewomen such as Clementina and Rusticiana.

At Constantinople, too, Gregory refused to surrender monastic community. A band
of brothers from Saint Andrew's---among them his close friend Maximian, later
bishop of Syracuse---followed him east and lived beside him in the palace
quarters. Gregory treasured their presence as his spiritual lifeline:

\begin{quote}
Which happened, I perceive, by Divine dispensation, in order that by their
example, as by an anchored cable, I might ever be kept fast to the tranquil
shore of prayer, whenever I should be tossed by the ceaseless waves of secular
affairs. For to their society I fled as to the bosom of the safest port from
the rolling swell, and from the waves of earthly occupation; and though that
office which withdrew me from the monastery had with the point of its
employments stabbed me to death as to my former tranquillity of life, yet in
their society, by means of the appeals of diligent reading, I was animated with
the yearnings of daily renewed compunction. (\work{Moralia in Iob}, Ep.\ Ded.\
1)
\end{quote}

In that same circle Gregory met Bishop Leander of Seville, who had come to
Constantinople on an embassy concerning the Catholic faith among the Arian
Visigoths of Spain. Their friendship endured across the Mediterranean for the
rest of their lives. Together with the monks from Saint Andrew's, Leander
pressed Gregory to expound the Book of Job, asking him not only to trace the
historical text and its allegorical sense but to turn each mystery into moral
instruction fortified by scriptural testimonies (\work{Moralia}, Ep.\ Ded.\ 1).
Out of those oral conferences in Constantinople grew the monumental
\work{Moralia in Iob}.

\subsection{The dispute with Patriarch Eutychius}

Gregory's years in Constantinople also drew him into a sharp theological
confrontation that he himself recorded in full in \work{Moralia in Iob} XIV.56.
Eutychius, the elderly Patriarch of Constantinople, had circulated a treatise
on the resurrection arguing that the risen human body would be impalpable,
``more subtle than the wind and air.'' Commenting on Job 19:26 (``And I shall be
again encompassed with my skin''), Gregory relates how he challenged the
patriarch directly with Luke 24:39---Christ inviting his doubting disciples to
touch his flesh and bones. When Eutychius answered that Christ had rendered his
body temporarily palpable to cure the disciples' unbelief before letting it
revert to an impalpable subtlety, Gregory pressed the logic of Romans 6:9:
``If then there was aught in the Body which was capable of being altered after
His resurrection, contrary to the truly spoken declaration of Paul, the Lord
after His resurrection returned into death'' (\work{Moralia} XIV.56.72).

Eutychius next appealed to 1 Corinthians 15:50 (``Flesh and blood cannot
inherit the Kingdom of God''). Gregory replied with a crisp distinction:
Scripture speaks of flesh in two senses---``according to nature'' (as in
Genesis 2:23 and John 1:14, ``The Word was made flesh'') and ``according to sin
or corruption'' (as in Genesis 6:3 and Romans 8:9). In the kingdom of heaven
there will be flesh according to nature, incorruptible and subtle through
spiritual power, yet truly palpable in its reality, lest ``he that dies is one
person, and he that rises again is another'' (\work{Moralia} XIV.56.72, 77).
When the dispute grew heated, Emperor Tiberius II summoned both men to a
private audience, weighed their arguments, judged Eutychius's treatise
erroneous, and ordered it burned. Both disputants fell gravely ill immediately
afterward. Gregory recovered from violent fever; Eutychius died (582), but not
before gripping the skin of his own hand in front of Gregory's visiting friends
and confessing, ``I confess that we shall all rise again in this flesh''
(\work{Moralia} XIV.56.74). Out of charity, Gregory refrained from pressing the
polemic further ``lest I should seem to be darting words at his ashes.''

Around 585 or 586, Pelagius II recalled Gregory to Rome. Returning gladly to
Saint Andrew's on the Caelian Hill---probably serving now as its abbot while
continuing to advise the pope, including drafting a major letter to the
schismatic bishops of Istria on the Three Chapters as Paul the Deacon later
reported (Barmby, Prolegomena)---Gregory enjoyed roughly five years of monastic
life. He would look back on them for the rest of his days as the happiest
season he ever knew.
